% !TeX root = ../sections/03_solutions.tex
\subsection{対策-シュミット-直交化}
\begin{background}
	シュミットの直交化法には，正規化を伴う場合と伴わない場合がある。
	
	正規化を伴わない場合は，
	\[
	(\varphi_i,\varphi_j)=0\quad (i\neq j)
	\]
	だけを満たす直交関数列を作る。
	
	正規化を伴う場合は，さらに各関数の長さを \(1\) にそろえる。
	つまり，
	\[
	(e_i,e_j)=0\quad (i\neq j),
	\qquad
	(e_i,e_i)=1
	\]
	を満たすようにする。
	
	正規化を伴わない直交関数 \(\varphi_i\) が得られたあと，
	\[
	e_i=\frac{\varphi_i}{\|\varphi_i\|}
	\]
	とすれば，正規化された直交関数が得られる。
\end{background}

\begin{strategy}
	今回得られた直交関数列は
	\[
	\varphi_1=x,\qquad
	\varphi_2=x^2,\qquad
	\varphi_3=3-5x^2
	\]
	である。
	
	それぞれのノルムを計算すると，
	\[
	\|\varphi_1\|^2
	=
	\int_{-1}^{1}x^2\,dx
	=
	\frac{2}{3},
	\]
	\[
	\|\varphi_2\|^2
	=
	\int_{-1}^{1}x^4\,dx
	=
	\frac{2}{5},
	\]
	\[
	\|\varphi_3\|^2
	=
	\int_{-1}^{1}(3-5x^2)^2\,dx
	=
	8.
	\]
	
	したがって，正規化するなら
	\[
	e_1=\sqrt{\frac{3}{2}}x,
	\]
	\[
	e_2=\sqrt{\frac{5}{2}}x^2,
	\]
	\[
	e_3=\frac{3-5x^2}{2\sqrt{2}}
	\]
	となる。
	
	ただし，問題文に「正規化を伴わない」とある場合は，
	\[
	\varphi_1=x,\qquad
	\varphi_2=x^2,\qquad
	\varphi_3=3-5x^2
	\]
	または
	\[
	\varphi_3=1-\frac{5}{3}x^2
	\]
	の形で答えればよい。
\end{strategy}